% GENERATED FILE. Edit canonical lessons, then run npm run generate:books.
% canonical-source-hash: fnv1a64:55a9c6ff559737d8
% canonical-lessons: IT-C451-lungomare, IT-C451-gironzolare, IT-C451-curiosare, IT-C451-sbirciare, IT-C451-vetrina

\chapter{Due passi in città}
\label{ch:it-due-passi-in-citta}
\hlchaptermodality{451}

\begin{chapteropening}
\textbf{By the end of this chapter:} \emph{I can talk about strolling round a town: the seafront, mooching about, browsing, peeking and the shop window.}

Complete the last lesson of chapter 451: I can talk about strolling round a town: the seafront, mooching about, browsing, peeking and the shop window.
\end{chapteropening}

\section[il lungomare]{il lungomare — the seafront, the promenade}

\label{lesson:IT-C451-lungomare}



\begin{quote}
Before the new one: say the Italian for the patron saint, then the Italian for the stall.
\end{quote}

\subsection*{You'll want to know: il lungomare}
\textbf{il lungomare} — \textquotedblleft{}the seafront, the promenade\textquotedblright{}.

\begin{quote}
La sera passeggiamo sul lungomare.
\end{quote}

\begin{quote}
\emph{In the evening we stroll on the seafront.}
\end{quote}

By a lake it is \emph{il lungolago}.

\begin{cousinweb}
\emph{lungomare} is built from \emph{lungo} 'along', inherited from Latin \emph{longus} 'long', and \emph{mare} 'sea', inherited from Latin \emph{mare} (EtymDB).
\end{cousinweb}

\begin{grammarlens}[title={What you've built}]
One more word for telling, explaining and describing at B1.
\end{grammarlens}

\subsection*{Your turn}
\begin{itemize}
\raggedright
  \item \emph{Say it:} \emph{il lungomare}
  \item \emph{Say it:} \emph{il lungomare}, once more
  \item \emph{Say it:} a sentence of your own with \emph{il lungomare}
  \item \emph{Recall:} say the Italian for the patron saint, then the Italian for the stall, then say \emph{il lungomare} again
\end{itemize}

\begin{tcolorbox}[breakable,colback=teal!4,colframe=teal!35!black,title={Before you move on}]
What does il lungomare mean? (\textquotedblleft{}The seafront, the promenade\textquotedblright{}.) Say it once more.
\end{tcolorbox}
\section[gironzolare]{gironzolare — to wander around, to mooch about}

\label{lesson:IT-C451-gironzolare}



\begin{quote}
Before the new one: say the Italian for the stall, then the Italian for the seafront.
\end{quote}

\subsection*{You'll want to know: gironzolare}
\textbf{gironzolare} — \textquotedblleft{}to wander around, to mooch about\textquotedblright{}.

\begin{quote}
Ho gironzolato per i negozi senza comprare niente.
\end{quote}

\begin{quote}
\emph{I mooched around the shops without buying anything.}
\end{quote}

Informal; built on \emph{girare}.

\begin{cousinweb}
\emph{gironzolare} is built on \emph{girare} 'to turn', which is inherited from Latin \emph{gyrāre}, from Greek \emph{gŷros} 'circle' (Lewis \& Short): going round in circles.
\end{cousinweb}

\begin{grammarlens}[title={What you've built}]
One more word for telling, explaining and describing at B1.
\end{grammarlens}

\subsection*{Your turn}
\begin{itemize}
\raggedright
  \item \emph{Say it:} \emph{gironzolare}
  \item \emph{Say it:} \emph{gironzolare}, once more
  \item \emph{Say it:} a sentence of your own with \emph{gironzolare}
  \item \emph{Recall:} say the Italian for the stall, then the Italian for the seafront, then say \emph{gironzolare} again
\end{itemize}

\begin{tcolorbox}[breakable,colback=teal!4,colframe=teal!35!black,title={Before you move on}]
What does gironzolare mean? (\textquotedblleft{}To wander around, to mooch about\textquotedblright{}.) Say it once more.
\end{tcolorbox}
\section[curiosare]{curiosare — to browse, to have a look round; to snoop}

\label{lesson:IT-C451-curiosare}



\begin{quote}
Before the new one: say the Italian for the seafront, then the Italian for to wander around.
\end{quote}

\subsection*{You'll want to know: curiosare}
\textbf{curiosare} — \textquotedblleft{}to browse, to have a look round; to snoop\textquotedblright{}.

\begin{quote}
Entriamo a curiosare in libreria.
\end{quote}

\begin{quote}
\emph{Let us go into the bookshop for a look round.}
\end{quote}

\emph{Curioso} — curious.

\begin{cousinweb}
\emph{curiosare} is built on \emph{curioso}, a learned borrowing from Latin \emph{cūriōsus} 'careful, inquisitive' (EtymDB), from \emph{cūra} 'care'.
\end{cousinweb}

\begin{grammarlens}[title={What you've built}]
One more word for telling, explaining and describing at B1.
\end{grammarlens}

\subsection*{Your turn}
\begin{itemize}
\raggedright
  \item \emph{Say it:} \emph{curiosare}
  \item \emph{Say it:} \emph{curiosare}, once more
  \item \emph{Say it:} a sentence of your own with \emph{curiosare}
  \item \emph{Recall:} say the Italian for the seafront, then the Italian for to wander around, then say \emph{curiosare} again
\end{itemize}

\begin{tcolorbox}[breakable,colback=teal!4,colframe=teal!35!black,title={Before you move on}]
What does curiosare mean? (\textquotedblleft{}To browse, to have a look round; to snoop\textquotedblright{}.) Say it once more.
\end{tcolorbox}
\section[sbirciare]{sbirciare — to peek, to sneak a look}

\label{lesson:IT-C451-sbirciare}



\begin{quote}
Before the new one: say the Italian for to wander around, then the Italian for to browse.
\end{quote}

\subsection*{You'll want to know: sbirciare}
\textbf{sbirciare} — \textquotedblleft{}to peek, to sneak a look\textquotedblright{}.

\begin{quote}
Ho sbirciato dalla finestra.
\end{quote}

\begin{quote}
\emph{I peeked out of the window.}
\end{quote}

Quick and secret.

\begin{cousinweb}
\emph{sbirciare} is usually linked to an old word \emph{bircio} 'short-sighted', but the origin is not settled.
\end{cousinweb}

\begin{grammarlens}[title={What you've built}]
One more word for telling, explaining and describing at B1.
\end{grammarlens}

\subsection*{Your turn}
\begin{itemize}
\raggedright
  \item \emph{Say it:} \emph{sbirciare}
  \item \emph{Say it:} \emph{sbirciare}, once more
  \item \emph{Say it:} a sentence of your own with \emph{sbirciare}
  \item \emph{Recall:} say the Italian for to wander around, then the Italian for to browse, then say \emph{sbirciare} again
\end{itemize}

\begin{tcolorbox}[breakable,colback=teal!4,colframe=teal!35!black,title={Before you move on}]
What does sbirciare mean? (\textquotedblleft{}To peek, to sneak a look\textquotedblright{}.) Say it once more.
\end{tcolorbox}
\section[la vetrina]{la vetrina — the shop window}

\label{lesson:IT-C451-vetrina}



\begin{quote}
Before the new one: say the Italian for to browse, then the Italian for to peek.
\end{quote}

\subsection*{You'll want to know: la vetrina}
\textbf{la vetrina} — \textquotedblleft{}the shop window\textquotedblright{}.

\begin{quote}
Ho visto delle scarpe in vetrina.
\end{quote}

\begin{quote}
\emph{I saw some shoes in the window.}
\end{quote}

\emph{Guardare le vetrine} — to go window-shopping.

\begin{cousinweb}
\emph{vetrina} is built on \emph{vetro} 'glass', which is inherited from Latin \emph{vitrum} (Lewis \& Short).
\end{cousinweb}

\begin{grammarlens}[title={What you've built}]
One more word for telling, explaining and describing at B1.
\end{grammarlens}

\subsection*{Your turn}
\begin{itemize}
\raggedright
  \item \emph{Say it:} \emph{la vetrina}
  \item \emph{Say it:} \emph{la vetrina}, once more
  \item \emph{Say it:} a sentence of your own with \emph{la vetrina}
  \item \emph{Recall:} say the Italian for to browse, then the Italian for to peek, then say \emph{la vetrina} again
\end{itemize}

\begin{tcolorbox}[breakable,colback=teal!4,colframe=teal!35!black,title={Before you move on}]
What does la vetrina mean? (\textquotedblleft{}The shop window\textquotedblright{}.) Say it once more.
\end{tcolorbox}
